\section{引戻し成分、対数型成分、および交叉の厚み}
\label{sec:allm-geometry-v33}

\subsection{重み付き射影平面への引戻し}

$X_m=\mathbb P_{\mathrm{stack}}(1,1,1,m)$ から
$Y_m=\mathbb P_{\mathrm{stack}}(1,1,m)$ への一般の重み付き射影を考える。
$Y_m$ の同次座標を $(x,y,w)$ とすれば、重み $m+3$ の水平一形式は一意に
\[
 \eta=c_3\,dw-\frac m3w\,dc_3
       +\left(f_{m+1}-\frac m3wL_1\right)(x\,dy-y\,dx)
\]
と書ける。$c_3,L_1,f_{m+1}$ は二変数の、それぞれ三次、一次、$m+1$ 次の式である。
係数空間の次元は
\[
 4+2+(m+2)=m+8
\]
となる。
この表示は、$dw$ 係数を $c_3$ と置き、水平性から $w$ を含む項を解き、
残りの二変数の水平形式が $(x\,dy-y\,dx)$ の倍数になることから従う。
$m>3$ では次数だけから $dw$ 係数が二変数三次式となる。
$m=3$ では一見 $\lambda w$ も $dw$ 係数に許されるが、水平性を表す式の
$w^2$ 係数は $3\lambda$ だけなので $\lambda=0$ である。
ここで $dx,dy$ の係数の次数は $m+2<2m$ であり、$w^2$ を含めないことを用いた。

二次元の標的では可積分条件は自動的に成立する。
従って、これらの形式を $X_m$ へ引き戻し、
$\wedge^2T_{X_m}\simeq\Omega^1_{X_m}\otimes K_{X_m}^{-1}$ により
対応する Poisson テンソルを取れば、$\mathscr P_m$ の族を得る。
この族の閉包を $\mathscr C_m$ と書く。

射影の自由度は、三次元の一次形式空間の二平面を選ぶ $2$ 方向と、
$w$ に加える次数 $m$ の式を二変数式で割った $N_m-(m+1)$ 方向である。
従って総パラメータ数は
\[
 2+\{N_m-(m+1)\}+(m+8)=N_m+9.
\]
これは一般ファイバーの次元を無視した数え上げではない。
一般の $c_3$ は三つの異なる共点直線に分解するため、選んだ二平面はその三因子の
張る空間から復元できる。
二平面を固定すれば、非二変数の $w$ 移動は、前節で求めた
重量 $0,\ldots,m-1$ の混合核を一意に与える。
標的形式による残りの自由度は重量 $-1$ にある。
従って、このパラメータ表示は一般点で余分な正次元ファイバーを持たない。

\begin{corollary}[任意に大きい一般点の長さを持つ成分]
\label{cor:allm-pullback-component-v33}
各 $m\geq3$ に対し、$\mathscr C_m$ は $\mathscr P_m$ の既約成分であり、
\[
 \dim\mathscr C_m=N_m+9.
\]
その一般点の局所環の長さは $m+3$ である。
したがって、重み付き射影三次元スタック上の大域的 Poisson スキームには、
一般点の非被約な長さが任意に大きくなる成分が存在する。
\end{corollary}

\begin{proof}
パラメータ空間は既約であり、上の復元から像の次元は $N_m+9$ である。
一般の二変数三次式は線形座標変換と倍率によって $xy(x+y)$ に直せる。
従って前節の局所模型における $a=b=0$ の枝を、この引戻し族が一般点で覆う。
その枝の次元も $N_m+9$ なので、閉包は一つの既約成分である。
この枝上の一般点では残余混合座標 $u$ が単元となる。
定理~\ref{thm:allm-concurrent-v33} の横断環から長さ $m+3$ を得る。
\end{proof}

全ての次元はアフィン係数空間での次元である。
Poisson テンソルの倍率や重み付き座標変換群で割った次元ではない。
この節は全ての既約成分の分類も、成分の最大次元の決定も主張しない。

\subsection{対数型の枝と二つの primary 成分}

前節の局所模型で $b=u=0$、$a\ne0$ とすると、$c$ は三つの非共点直線の積となる。
混合項は次数 $m$ の $w$ 移動で消去でき、三角形座標では基本一形式が
\[
 XYZ w\left(\lambda_0\frac{dw}{w}
 +\lambda_1\frac{dX}{X}
 +\lambda_2\frac{dY}{Y}
 +\lambda_3\frac{dZ}{Z}\right),
 \qquad m\lambda_0+\lambda_1+\lambda_2+\lambda_3=0
\]
となる。従ってこの枝は対数型 $L(1,1,1,m)$ の閉包の局所枝である。
逆に、重み $m$ の因子の $w$ 係数が非零なら、その因子は $w+g_m$ と書けるので、
一般の対数型テンソルもこの表示に含まれる。
この対数型成分の次元も $N_m+9$ であり、一般点では被約である。

滑らかな因子を省き、$n=m+2$ と置く。局所模型の関係イデアルには
正確な primary 分解
\[
 (ab,b^2,a^n u)=(a^n,ab,b^2)\cap(b,u)
\]
がある。
第一のイデアルは $(a,b)$-primary、第二は素イデアルであり、embedded prime はない。
二つの被約枝の交叉は滑らかだが、引戻し側の primary 成分と対数型枝の
スキーム論的交叉は
\[
 k[[a,b,u]]/(a^n,b,u)\simeq k[a]/(a^n)
\]
である。従って、滑らかな被約交叉の上に長さ $m+2$ の厚みを持つ。

さらに Hilbert--Burch 行列
\[
 \begin{pmatrix}
 b&a^{n-1}u\\
 -a&0\\
 0&-b
 \end{pmatrix}
\]
の最大小行列式は、符号を除いて $ab,b^2,a^nu$ である。
関係イデアルの高さは二なので、滑らかな因子を除いた一次元局所環は
Cohen--Macaulay である。一方、三つの生成元は最小なので完全交叉ではない。
一般の横断超平面を $u=-a-\beta b$ と取ると、切断環は
\[
 k[a,b]/(a^{n+1},ab,b^2)
\]
となり、長さは $n+2=m+4$ である。
これは二枝の一般点での長さ $(m+3)+1$ とも一致する。

\begin{remark}[原点の形式変形問題との接続]
ここでの非被約な長さは、Poisson スキームの非零点で計算されたものである。
原点で完備化した形式変形基底に対する帰結には、斉次錐の既約成分とその一般点の
長さが原点完備化により保たれることを別に用いる。
その議論と外側の例外列代数の二次 miniversal hull の同定を組み合わせることで、
同じ長さ $m+3$ が圏の変形基底にも現れる。
この接続は、任意の Poisson 量子化が AS 正則であるという主張を必要としない。
\end{remark}
